\documentclass[12pt]{article}
\usepackage{latexblog}

% ============================================================
% Blog Metadata
% ============================================================
\blogtitle{Java9模块(Modules)}
\blogdate{2022-05-10}
\blogtags{编程语言, 闲话编程}
\bloglang{zh}
\blogsummary{}

\begin{document}
\makeblogtitle

Java9中引入的最重要的特性是模块化------Java Platform Module System (JPMS)。实现模块化是一件非常有挑战性的事情，在此之前，已经有过一些非官方的解决方案例如OSGI。实际在2005年Java 7的时候\href{https://jcp.org/en/jsr/detail?id=277}{JSR 277: Java Module System}就提议要模块化Java SE，后面该提议又被\href{https://jcp.org/en/jsr/detail?id=376}{JSR 376: Java Platform Module System}取代并推延到Java 8，又推延到Java 9，最终直到2017年Java 9推迟后才发布出来。

\section{目标}\label{ux76eeux6807}

JSR 376中阐述了模块化Java SE平台的关键目标：

\begin{itemize}
\tightlist
\item
  可靠的依赖配置： 必须要能够通过模块化系统显式定义模块的依赖信息，能够在编译期和运行时都能够有效识别，并可以通过遍历所有模块的方式得到程序运行所需的所有模块信息。
\item
  高度封装性：只有当模块中的包显式暴露、且被其他模块显式依赖的时候才能使用。
\item
  对Java 平台进行定制：可以对特定平台删除或者添加模块，从而实现Java平台的裁剪。
\item
  隐藏Java平台细节：一些不希望暴露给用户的方法可以通过模块化隐藏起来。
\item
  提高运行效率：在预知模块依赖的情况下，JVM有一些优化可以提高性能。
\end{itemize}

\section{模块化的使用}\label{ux6a21ux5757ux5316ux7684ux4f7fux7528}

\subsection{JDK的模块}\label{jdkux7684ux6a21ux5757}

JDK本身也拆分成了一系列的模块，可以通过命令列出来:

\begin{Shaded}
\begin{Highlighting}[]
\ExtensionTok{java} \AttributeTok{{-}{-}list{-}modules}                                   

\ExtensionTok{java.base@11.0.13}
\ExtensionTok{java.compiler@11.0.13}
\ExtensionTok{java.datatransfer@11.0.13}
\ExtensionTok{java.desktop@11.0.13}
\ExtensionTok{java.instrument@11.0.13}
\ExtensionTok{java.logging@11.0.13}
\ExtensionTok{java.management@11.0.13}
\ExtensionTok{...}
\end{Highlighting}
\end{Shaded}

\subsection{如何创建一个模块}\label{ux5982ux4f55ux521bux5efaux4e00ux4e2aux6a21ux5757}

创建模块十分简单，只需要添加一个特殊的module-info.java文件，并添加模块的描述信息(requires, exports, provides..with, uses, opens)：

\begin{Shaded}
\begin{Highlighting}[]
\NormalTok{module modulename }\OperatorTok{\{}
\NormalTok{    requires module}\OperatorTok{.}\FunctionTok{name}\OperatorTok{;}
\OperatorTok{\}}
\end{Highlighting}
\end{Shaded}

\begin{itemize}
\tightlist
\item
  \texttt{requires}用来申明模块依赖于其他模块
\item
  \texttt{requires\ static}用来申明编译时依赖（运行时是可选的）
\item
  \texttt{requires\ transtive}表示如果A依赖某个模块B，那么依赖A的模块也隐式依赖B
\item
  \texttt{exports}导出包下的public类型（包括嵌套的public和protected类型）
\item
  \texttt{exports...to}仅导出给指定的模块
\item
  \texttt{uses}申明该模块为服务的消费者（使用了某接口或者抽象类）
\item
  \texttt{provides...with}申明模块为服务ud提供者（提供了实现）
\item
  \texttt{opens}申明包对于其他模块在运行时（包括反射）可见
\item
  \texttt{opens..to}申明包对于特定的模块模块在运行时（包括反射）可见
\item
  \texttt{open}申明模块下的所有包均在运行时可见
\end{itemize}

\begin{Shaded}
\begin{Highlighting}[]
\NormalTok{module my}\OperatorTok{.}\FunctionTok{module} \OperatorTok{\{}
\NormalTok{    exports com}\OperatorTok{.}\FunctionTok{my}\OperatorTok{.}\FunctionTok{package}\OperatorTok{.}\FunctionTok{name}\OperatorTok{;}
\OperatorTok{\}}

\NormalTok{module my}\OperatorTok{.}\FunctionTok{module} \OperatorTok{\{}
\NormalTok{    export com}\OperatorTok{.}\FunctionTok{my}\OperatorTok{.}\FunctionTok{package}\OperatorTok{.}\FunctionTok{name}\NormalTok{ to com}\OperatorTok{.}\FunctionTok{specific}\OperatorTok{.}\FunctionTok{package}\OperatorTok{;}
\OperatorTok{\}}
\end{Highlighting}
\end{Shaded}

\begin{Shaded}
\begin{Highlighting}[]
\NormalTok{module my}\OperatorTok{.}\FunctionTok{module} \OperatorTok{\{}
\NormalTok{    uses }\KeywordTok{class}\OperatorTok{.}\FunctionTok{name}\OperatorTok{;}
\OperatorTok{\}}

\NormalTok{module my}\OperatorTok{.}\FunctionTok{module} \OperatorTok{\{}
\NormalTok{    provides MyInterface with MyInterfaceImpl}\OperatorTok{;}
\OperatorTok{\}}
\end{Highlighting}
\end{Shaded}

默认情况下，private在模块中是不可见的（即使通过反射的方式），因此如果想暴露出来，可以通过\texttt{open}来实现。

\begin{Shaded}
\begin{Highlighting}[]
\NormalTok{module foo }\OperatorTok{\{}
\NormalTok{    opens com}\OperatorTok{.}\FunctionTok{example}\OperatorTok{.}\FunctionTok{bar}\OperatorTok{;}
\OperatorTok{\}}

\NormalTok{module my}\OperatorTok{.}\FunctionTok{module} \OperatorTok{\{}
\NormalTok{    opens com}\OperatorTok{.}\FunctionTok{my}\OperatorTok{.}\FunctionTok{package}\NormalTok{ to moduleOne}\OperatorTok{,}\NormalTok{ moduleTwo}\OperatorTok{,}\NormalTok{ etc}\OperatorTok{.;}
\OperatorTok{\}}

\NormalTok{open module foo }\OperatorTok{\{}
\OperatorTok{\}}
\end{Highlighting}
\end{Shaded}

\section{Example：创建一个Java FX程序}\label{exampleux521bux5efaux4e00ux4e2ajava-fxux7a0bux5e8f}

首先，需要通过gradle插件来自动引入JavaFX的依赖项:

\begin{Shaded}
\begin{Highlighting}[]
\NormalTok{plugins }\OperatorTok{\{}
\NormalTok{    id }\StringTok{\textquotesingle{}java\textquotesingle{}}
\NormalTok{    id }\StringTok{\textquotesingle{}application\textquotesingle{}}
\NormalTok{    id }\StringTok{\textquotesingle{}org.openjfx.javafxplugin\textquotesingle{}}\NormalTok{ version }\StringTok{\textquotesingle{}0.0.10\textquotesingle{}}
\OperatorTok{\}}

\NormalTok{javafx }\OperatorTok{\{}
\NormalTok{    version }\OperatorTok{=} \StringTok{"17.0.1"}
\NormalTok{    modules }\OperatorTok{=} \OperatorTok{[} \StringTok{\textquotesingle{}javafx.controls\textquotesingle{}} \OperatorTok{]}
\OperatorTok{\}}
\end{Highlighting}
\end{Shaded}

然后，在module-info.java中申明依赖关系:

\begin{Shaded}
\begin{Highlighting}[]
\NormalTok{module pulsar}\OperatorTok{.}\FunctionTok{browser} \OperatorTok{\{}
\NormalTok{    requires javafx}\OperatorTok{.}\FunctionTok{controls}\OperatorTok{;}
\NormalTok{    exports com}\OperatorTok{.}\FunctionTok{riguz}\OperatorTok{.}\FunctionTok{pulsar}\OperatorTok{.}\FunctionTok{browser}\OperatorTok{;}
\OperatorTok{\}}
\end{Highlighting}
\end{Shaded}

Ref:

\begin{itemize}
\tightlist
\item
  \href{https://www.oracle.com/hk/corporate/features/understanding-java-9-modules.html}{Understanding Java 9 Modules}
\item
  \href{https://openjdk.java.net/jeps/200}{JEP 200: The Modular JDK}
\end{itemize}


\end{document}
